% ============================================================================
%  preamble/code.tex  --  syntax-highlighted code on slides
% ----------------------------------------------------------------------------
%  Uses `listings`, which is pure LaTeX: no Python, no --shell-escape, works
%  on Overleaf and on a fresh machine with nothing installed.
%
%  IMPORTANT: any frame containing code must be marked fragile:
%      \begin{frame}[fragile]{Title}
%        \begin{lstlisting}[language=Python]
%        ...
%        \end{lstlisting}
%      \end{frame}
%  Without [fragile] LaTeX will fail with a confusing error. This is a beamer
%  quirk, not a bug in this template.
% ============================================================================

\usepackage{listings}

% Colours are taken from the theme palette so code matches the slide design.
\lstdefinestyle{deck}{
  basicstyle       = \ttfamily\scriptsize\color{BrandInk},
  keywordstyle     = \color{BrandPrimary}\bfseries,
  commentstyle     = \color{BrandMuted}\itshape,
  stringstyle      = \color{BrandAccent},
  numberstyle      = \ttfamily\tiny\color{BrandMuted},
  identifierstyle  = \color{BrandInk},
  backgroundcolor  = \color{BrandWash},
  %
  numbers          = left,
  numbersep        = 9pt,
  stepnumber       = 1,
  %
  %  No frame rule. `frame=leftline` puts a rule per line, which fragments
  %  visually once line numbers sit beside it. The background wash already
  %  separates code from prose, which is all a slide needs.
  frame            = none,
  xleftmargin      = 22pt,
  xrightmargin     = 4pt,
  framexleftmargin = 22pt,
  %
  breaklines       = true,
  breakatwhitespace= true,
  breakindent      = 1.5em,
  breakautoindent  = true,
  %
  %  NOTE: `postbreak` (the little arrow marking a wrapped line) is
  %  deliberately NOT set. It clashes with the way listings kerns lines
  %  inside a beamer [fragile] frame and produces a flood of
  %  "Missing number, treated as zero" errors. Line breaking itself works
  %  fine without it.
  %
  showstringspaces = false,
  tabsize          = 2,
  columns          = flexible,
  keepspaces       = true,
  upquote          = true,
  aboveskip        = 1em,
  belowskip        = 0.6em,
}

\lstdefinelanguage{SPARQL}{
  morekeywords={PREFIX,BASE,SELECT,DISTINCT,REDUCED,CONSTRUCT,ASK,DESCRIBE,%
    WHERE,FROM,NAMED,GRAPH,OPTIONAL,UNION,MINUS,FILTER,BIND,VALUES,SERVICE,%
    ORDER,BY,ASC,DESC,GROUP,HAVING,LIMIT,OFFSET,AS,a},
  morekeywords=[2]{STR,LANG,DATATYPE,BOUND,IRI,REGEX,CONTAINS,STRSTARTS,%
    STRENDS,COUNT,SUM,MIN,MAX,AVG,COALESCE,IF},
  sensitive     = true,      % keep keywords UPPERCASE; lets lone `a` match too
  morecomment   = [l]{\#},   % SPARQL comments are #, not --
  morestring    = [b]",      % "literals"
  morestring    = [b]',      % 'literals'
}

\lstset{style=deck}

% ---------------------------------------------------------------------------
%  Inline code: \code{foo()} -- reads better than \texttt{} on slides
% ---------------------------------------------------------------------------
\newcommand{\code}[1]{%
  \colorbox{BrandWash}{\ttfamily\small\color{BrandPrimary}#1}%
}
